\section{Introduction}
\label{sec:intro}
Urban traffic-management centres typically observe congestion as it happens. Linking an \emph{advance} forecast to an
explainable, evidence-grounded and tested operational recommendation remains difficult: traffic depends on recurrent
demand, neighbouring links, weather, planned events, disruptions and infrastructure condition, and a forecast alone does
not establish which intervention helps. An LLM-generated signal recommendation is, by itself, not evidence that an
action is safe or effective.

We study a decision-support pipeline for the scenario of a permitted event near a busy corridor in which rain and a
nearby collision coincide. The pipeline (i) forecasts link speeds 15, 30 and 60 minutes ahead with a context-gated
graph spatiotemporal Transformer, (ii) attaches split-conformal prediction intervals, (iii) explains the forecast of
the actual model with grouped Shapley values, (iv) retrieves operational guidance with a FAISS index, (v) asks a local
LLM (Qwen2.5-7B-Instruct via Ollama) for schema-constrained candidate actions, (vi) applies deterministic safety
constraints, and (vii) tests admissible candidates in a SUMO microscopic simulation before human review.

\textbf{Proposed contributions} (stated as proposals; their effects are reported only where measured):
(1) availability-aware context fusion in the graph forecaster; (2) an explicit uncertainty-and-evidence gate that
decides whether advisory generation proceeds; (3) a simulation-based candidate gate with side-road protection; and
(4) a traceable evaluation that links forecast quality, grounding, deferral and simulated intervention outcomes to
versioned artifacts. We do not claim deployment, real-world congestion reduction or causal explanations.
